\clearpage
\fancyhead[L]{\small PIAZZI. LEZIONI DI ASTRONOMIA.}
\hypertarget{native-622-lower}{}
\begin{center}\small Göttingische gelehrte Anzeigen. Stück 4. Seiten 33..37. 6. Januar 1820.\\[3pt]\rule{.72\linewidth}{.3pt}\end{center}

\noindent\emph{Lezioni elementari di astronomia ad uso del real osservatorio di Palermo. Palermo {\upshape 1817}. Dalla stamperia reale.} Band I, 240 Seiten. Band II, 446 Seiten. Kleinquart.

\begingroup\brokenpenalty=10000\relax
Der Verfasser dieses Lehrbuchs, der berühmte und verdiente \textsc{Piazzi}, erklärt in der Vorrede, das Buch sei zunächst aus den Vorlesungen hervorgegangen, die er an der Sternwarte in Palermo gehalten habe, und sei eigentlich nicht für die Öffentlichkeit bestimmt gewesen. Auf den dringenden Wunsch der Schüler hin hätten die für das öffentliche Unterrichtswesen zuständigen Stellen die Veröffentlichung angeordnet. Von Schriften dieser Art erwartet man keine eigentliche Erweiterung der Wissenschaft selbst. Das Wichtigste ist eine klare, gründliche und ihrem Zweck angemessene Gliederung der Darstellung. Die Anforderungen an die Darstellung unterscheiden sich wiederum stark je nach Vorkenntnissen und Absichten der Schüler. Für bloße Dilettanten gelten andere Anforderungen als für mathematisch Gebildete. Wieder andere gelten für jemanden, der Sinn dafür hat, in der Wissenschaft eines der vollkommensten und edelsten Produkte des menschlichen \emph{Geistes} zu verehren, der nichts auf Treu und Glauben annimmt, sondern überall einen streng logischen Zusammenhang verlangt, und dem die vollendetste Darstellung am liebsten ist, selbst wenn sie mehr als gewöhnliche Anstrengung und Hilfskenntnisse erfordert. Schließlich gelten wieder andere Anforderungen für jemanden, der nur das niedrigere Ziel verfolgt, sich die Fertigkeit anzueignen, Himmelserscheinungen zu beobachten und zu berechnen, soweit es eine beschränktere mathematische Vorbildung erlaubt. Der Verfasser gibt zwar nicht ausdrücklich an, für welche Gruppe von Schülern er dieses Buch eigentlich bestimmt, doch scheint er vor allem an die letzte Gruppe gedacht zu haben. Von einem Astronomen wie \textsc{Piazzi}, der sich um mehrere Teile der Astronomie so große Verdienste erworben hat, war im Übrigen zu erwarten, dass auch ein elementares Werk dieser Art viele eigenständige Ansichten, bemerkenswerte Äußerungen und Urteile enthalten musste. In dieser Hinsicht ist es selbst für Astronomen nicht ohne Interesse.

Der Verfasser hat sein Werk in sieben Bücher eingeteilt. Das erste gibt einen allgemeinen Überblick über die Himmelserscheinungen, soweit man ihn allein durch grobe Beobachtung, ohne Instrumente und ohne besondere mathematische Kenntnisse gewinnen kann. Behandelt werden hier also die allgemeine tägliche Bewegung, die auf der Himmelskugel angenommenen Punkte und Kreise, die Gestalt der Erde, der Unterschied zwischen Fixsternen und Planeten, die Sonne, der Mond, die Planeten und Kometen, die Ekliptik und die übrigen dadurch bestimmten Kreise der Himmelskugel, der Tierkreis und die Sternbilder, das Vorrücken der Tagundnachtgleichen, die Finsternisse sowie die ersten Ungleichheiten der Planeten und ihre Erklärung durch das kopernikanische Weltsystem. Das zweite Buch behandelt die Grundlagen, die für die neue Astronomie vorausgesetzt werden. Hier finden wir einiges über die Lösung sphärischer Dreiecke und einen Überblick über die wichtigsten heute gebräuchlichen astronomischen Instrumente, etwa Fernrohre, Lot, Libelle, Verniers, Mikrometer, Mittagsfernrohr, Uhr, Vollkreis, Quadrant, Zenitsektor, Äquatoreal und Vervielfältigungskreise (über Letztere urteilt \textsc{Piazzi} ungünstig, doch scheint er das, was zu ihrem \SourceLeaf{623}Nachteil gesagt werden kann und auch von anderen gesagt worden ist, etwas übertrieben zu haben. Hätte er selbst, sagt \textsc{Piazzi}, mit einem Instrument dieser Art beobachten sollen, so hätte er \emph{nicht den hundertsten Teil} dessen geleistet, was er getan hat). Anschließend werden Refraktion, Parallaxe, Aberration und Nutation behandelt, soweit die Beobachtungen deswegen einer Korrektur bedürfen. Danach geht es um die Zeit und die Methoden der Zeitbestimmung und schließlich um die Beobachtungen, die zur Bestimmung der Polhöhe, der Schiefe der Ekliptik und der Stellung der Gestirne zum Äquator und zur Ekliptik dienen. Das dritte Buch befasst sich dann ausschließlich mit den Fixsternen. Sehr ausführlich behandelt es die Berechnung der Präzession, Nutation und Aberration, dann die jährliche Parallaxe der Fixsterne, ihre Eigenbewegung, die Sternverzeichnisse und die Besonderheiten der Fixsterne. Einen sonderbaren Fehlgriff bei der Bestimmung der Größe der jährlichen Präzession können wir hier nicht übergehen. \textsc{Piazzi} zieht von der beobachteten jährlichen Änderung der Länge von 15 Fixsternen, die sich aus dem Vergleich von \textsc{Bradleys} Bestimmung für 1755 mit \textsc{Piazzis} eigener für 1805 ergibt, sowohl die Wirkung der Verschiebung der Ekliptik als auch die der \emph{Eigenbewegung} ab und nimmt dann aus allen (übrigens sehr unterschiedlich ausfallenden) Ergebnissen das Mittel. Dies ist ein offenkundiger logischer Zirkelschluss, denn die Eigenbewegung der einzelnen Sterne lässt sich nur erkennen, indem man die beobachtete Bewegung mit der bereits als bekannt vorausgesetzten Präzession vergleicht. Bei einer folgerichtigen Rechnung hätten alle Sterne ein und dasselbe Ergebnis liefern sollen, und zwar genau das, von dem man ausgegangen war, um die Eigenbewegungen zu finden.

Gegenstand des vierten Buches ist die elliptische Bewegung der Planeten, größtenteils aus der Theoria motus corporum coelestium entlehnt, jedoch unter Weglassung aller Beweise.

Im fünften Buch werden die Himmelskörper unseres Sonnensystems einzeln betrachtet. Bemerkenswert ist, wie sich \textsc{Piazzi} auf S. 182 über die Bestimmung der Rotationszeit der Venus durch den verstorbenen \textsc{Schröter} äußert. `Niente di meno non manca, chi non poco dubiti di tali scoperte, e le consideri come non molto dissimili dall' altra del satellite, che \textsc{Short} e qualche altro credeva di aver veduto intorno a quaesta pianeta. Una illusione ottica, un fenomeno non bene osservato, accompagnati da un poco di immaginazione, hanno talora fatto avanzare delle congetture, e sognare delle esseri, che ben presto son poi caduti in piena dimenticanza (\textsc{Piazzi} scheint hier doch vergessen zu haben, dass er selbst auf S. 46 des ersten Bandes die Existenz eines Venustrabanten für sehr wahrscheinlich erklärt hatte, an den jetzt kaum noch ein anderer Astronom glaubt). Ebenso wendet er sich auf S. 230 gegen die brennenden Mondvulkane, die einige Astronomen angeblich gesehen haben.

Im sechsten Buch geht es um Finsternisse, Sternbedeckungen und die Durchgänge der unteren Planeten vor der Sonnenscheibe, im siebten schließlich um die Kometen. Zur Bestimmung ihrer Bahnen wird hier nur die indirekte Methode in derselben höchst unvollkommenen und, so kann man sagen, planlosen Gestalt angegeben, die die Astronomen zu ihrer Verwunderung in einem bekannten, vor einigen Jahren erschienenen Lehrbuch der Astronomie gefunden haben.
\par\clearpage\endgroup
